% =====================================================================
\section{6.1 정규화: 편향-분산 복습과 L1/L2 페널티}
\begin{frame}{손실함수에 페널티 항을 더한다}
  Week 4.1에서 봤듯 모델이 너무 유연하면(파라미터가 크고 자유로우면)
  분산이 커져 \textbf{과적합}한다. \kb{정규화(regularization)}는 손실함수에
  ``가중치가 너무 커지지 않도록'' 페널티 항을 더해 유효 유연성을
  인위적으로 낮추는 기법이다:
  \[
  J(w) = \underbrace{\frac{1}{2m}\sum_{i=1}^m (h_w(x^{(i)})-y^{(i)})^2}_{\text{원래 손실(적합도)}} +
  \underbrace{\lambda R(w)}_{\text{정규화 항(단순함을 요구)}}
  \]
  \begin{itemize}
    \item $\lambda$가 \textbf{0}이면 원래 손실과 같고(정규화 없음),
      커질수록 ``데이터에 딱 맞추기''보다 ``가중치를 작게''를 더 중요시.
    \item 4.1의 언어로: $\lambda$를 올리면 \textbf{분산은 줄고 편향은 늘어난다.}
  \end{itemize}
  \vskip0.4em
  \textbf{이 절의 순서}: (1) ``가중치 작음 = 단순함''이 왜 성립,
  (2) L2·L1이 구체적으로 무슨 일, (3) L1이 정확히 0을 만드는 대수·기하
  이유, (4) $\lambda$가 편향-분산 위의 ``손잡이''인지 숫자로 확인.
\end{frame}
